% =====================================================================
\chapter{TRIỂN KHAI VÀ LẬP TRÌNH HỆ THỐNG}\label{chap:trienkhai}
% =====================================================================

Chương này trình bày cách tổ chức mã nguồn (\S\ref{sec:module}), giao thức giữa các bo
(\S\ref{sec:giaothuc}), cách hiện thực và cấu hình từng thành phần trên vi điều khiển
(\S\ref{sec:chitiet}) và các bước triển khai lên phần cứng thực tế (\S\ref{chap:huongdanchay}).

\section{Các module mã nguồn của máy chủ}\label{sec:module}
\begin{table}[H]
\centering\small
\caption{Module firmware máy chủ (\file{Server/main/})}\label{tab:modules}
\begin{tabularx}{\linewidth}{@{}p{4.3cm}rX@{}}
\toprule
\textbf{Tệp} & \textbf{Dòng} & \textbf{Trách nhiệm} \\ \midrule
\file{main.c} & 43 & Khởi tạo NVS, khôi phục phiên, màn hình, còi, servo, mạng theo thứ tự \\
\file{parking/parking\_state.c} & 65 & Trạng thái bãi dùng chung, đồng hồ bãi (critical section) \\
\file{parking/json\_scan.c} & 81 & Hàm đọc JSON dùng chung, không cấp phát động \\
\file{parking/parking\_update.c} & 108 & Kiểm tra bản tin \code{parking\_update} và \code{climate} \\
\file{parking/rfid\_session.c} & 519 & Bản tin RFID, phiên gửi xe, phí, hóa đơn, thông báo, nhật ký, chống lặp, đối chiếu IR, lưu/khôi phục, đặt giờ, lượt thanh toán QR \\
\file{parking/payment.c} & 131 & Chuỗi VietQR (EMVCo, CRC16), tìm giao dịch khớp mã trong JSON của SePay \\
\file{parking/parking\_persist.c} & 47 & Lưu/đọc phiên, nhật ký, đồng hồ trong NVS \\
\file{barrier/servo.c} & 230 & PWM servo, hẹn giờ tự đóng, chế độ cháy \\
\file{buzzer/buzzer.c} & 29 & Còi tích cực mức thấp \\
\file{display/lcd\_display.c} & 706 & Giao diện TFT, màn hình thanh toán QR, màn hình cháy, ảnh động \\
\file{display/*\_gif.c} & dữ liệu & Khung ảnh động đã nén LZSS \\
\file{network/wifi.c} & 137 & Kết nối vào Wi-Fi chung (trạm), tự kết nối lại \\
\file{network/network\_service.c} & 87 & Khởi động mạng, mDNS, cập nhật IP cho mDNS \\
\file{network/udp/udp\_server.c} & 188 & Nhận và điều phối bản tin UDP, gửi ACK \\
\file{network/sepay\_client.c} & 127 & Tác vụ hỏi vòng API SePay khi có xe đang thanh toán \\
\file{network/http/\allowbreak http\_server.c} & 283 & 13 route HTTP, API JSON, xác nhận thanh toán, kiểm tra mã PIN \\
\bottomrule
\end{tabularx}
\end{table}
\section{Giao thức và giao diện lập trình}\label{sec:giaothuc}
\begin{lstlisting}[language=json,caption={Bản tin \code{parking\_update} (Sensor Client → máy chủ)}]
{"type":"parking_update","slots":[1,0,0,1,0],"water_mm":0.2,"fire":false}
\end{lstlisting}
\begin{lstlisting}[language=json,caption={Bản tin \code{rfid\_event} (RFID Client → máy chủ)}]
{"type":"rfid_event","card_id":"A1B2C3D4","gate":"in","temperature_c":28.0}
\end{lstlisting}
\begin{lstlisting}[language=json,caption={Bản tin \code{climate} (RFID Client → máy chủ, mỗi 5\,s)}]
{"type":"climate","temperature_c":28.0}
\end{lstlisting}

\paragraph{Xác nhận sự kiện quẹt thẻ.} Bản tin \code{rfid\_event} có thêm trường \code{seq}
(số nguyên 32 bit, khác 0) do RFID Client đặt; máy chủ trả lời bằng một gói JSON cùng
\code{seq}:
\begin{lstlisting}[language=json,caption={Sự kiện quẹt thẻ và gói xác nhận}]
{"type":"rfid_event","card_id":"A1B2C3D4","gate":"in","temperature_c":28.0,"seq":305441741}
{"type":"rfid_ack","seq":305441741}
\end{lstlisting}
\begin{itemize}[nosep]
  \item RFID Client gửi lại đúng gói đó mỗi 300\,ms (tối đa 6 lần) cho tới khi nhận được ACK.
  \item Máy chủ nhớ 16 giá trị \code{seq} gần nhất: gặp lại một \code{seq} đã xử lý là bản
        gửi lại, chỉ trả ACK chứ không mở cổng lần hai.
  \item \code{seq} bắt đầu từ một số ngẫu nhiên (bộ sinh số ngẫu nhiên phần cứng của
        ESP8266) mỗi lần khởi động rồi tăng dần, nên gần như không trùng với số máy chủ
        vừa thấy; client không cần lưu gì vào flash.
\end{itemize}

\begin{lydo}[Tại sao không ký (xác thực) gói UDP]
Ba bo nằm trong mạng Wi-Fi riêng của nhóm, được bảo vệ bằng WPA2; chỉ thiết bị
biết mật khẩu Wi-Fi mới gửi được gói tin. Ở quy mô mô hình, nhóm ưu tiên giao thức đơn
giản, đọc được trực tiếp trong log và dễ giả lập bằng script. Các thao tác nhạy cảm từ
trình duyệt (mở cổng, đặt giờ) vẫn cần mã PIN. Ký gói bằng mã xác thực (HMAC) được để
lại cho hướng phát triển.
\end{lydo}

\begin{table}[H]
\centering\small
\caption{Đặc tả trường dữ liệu và ràng buộc kiểm tra trên máy chủ}\label{tab:fields}
\begin{tabularx}{\linewidth}{@{}llX@{}}
\toprule
\textbf{Trường} & \textbf{Kiểu} & \textbf{Ràng buộc} \\ \midrule
\code{type} & chuỗi & \code{"parking\_update"}, \code{"rfid\_event"} hoặc \code{"climate"} \\
\code{seq} & số nguyên & Chỉ trong \code{rfid\_event}; $1 \le x \le 2^{32}-1$, không dấu, không phần thập phân \\
\code{slots} & mảng & Đúng 5 phần tử, mỗi phần tử 0 (trống) hoặc 1 (có xe) \\
\code{water\_mm} & số & $0 \le x \le 10\,000$ \\
\code{fire} & boolean & Tùy chọn (client cũ có thể không gửi); không chấp nhận 0/1 \\
\code{card\_id} & chuỗi & 1–31 ký tự \code{[A-Za-z0-9\_-]}, chuẩn hóa thành chữ hoa \\
\code{gate} & chuỗi & \code{"in"} hoặc \code{"out"} (không phân biệt hoa thường) \\
\code{temperature\_c} & số & $-40 \le x \le 125$ \\
Chung & — & Khóa lạ, khóa lặp, dữ liệu thừa sau \code{\}}, chuỗi có ký tự thoát đều bị từ chối \\
\bottomrule
\end{tabularx}
\end{table}

\begin{table}[H]
\centering\small
\caption[API HTTP của máy chủ (13 route)]{API HTTP của máy chủ (13 route); mọi lệnh POST cần header \code{X-Parking-Key} đúng mã PIN, sai thì trả 401}\label{tab:api}
\begin{tabularx}{\linewidth}{@{}lp{4.4cm}X@{}}
\toprule
\textbf{PT} & \textbf{Đường dẫn} & \textbf{Kết quả} \\ \midrule
GET & \code{/}, \code{/style.css} & Trang web nhúng trong firmware \\
GET & \code{/api/parking} & \code{received}, \code{online}, \code{slots}, \code{free}, \code{total}, \code{temperature\_c}, \code{water\_mm}, \code{fire}, \code{flood} \\
GET & \code{/api/barrier} & \code{\{"in\_open":bool,"out\_open":bool\}} \\
GET & \code{/api/rfid} & \code{now\_s}, \code{fee\_per\_hour}, \code{today}/\code{week} (doanh thu, lượt vào, lượt ra), \code{parked[]} (thẻ, giờ vào, phí tạm tính), \code{payment} (lượt thanh toán đang chờ: \code{\{"card","fee","code","expired"\}} hoặc \code{null}), \code{log[]} (10 sự kiện mới nhất) \\
POST & \code{/api/barrier/}\allowbreak\code{<cổng>/<lệnh>} & Cổng \code{in}, \code{out}, \code{all}; lệnh \code{open}, \code{close}. Lỗi 500 nếu bị từ chối (ví dụ đóng khi có cháy) \\
POST & \code{/api/time} & Thân là giờ Unix (giây) từ trình duyệt; máy chủ đặt đồng hồ bãi và lưu NVS \\
POST & \code{/api/payment/\allowbreak confirm} & Thân là mã thẻ; người vận hành xác nhận đã nhận tiền (chuyển khoản hoặc tiền mặt): xe ra như bình thường. Dùng được cả cho xe chưa quẹt ra (thu phí đến hiện tại) \\
\bottomrule
\end{tabularx}
\end{table}

\section{Hiện thực từng thành phần}\label{sec:chitiet}

\subsection{Thanh chắn: PWM phần cứng và hẹn giờ một lần}
Servo SG90 nhận xung chu kỳ 20\,ms; độ rộng xung quyết định góc (Hình~\ref{fig:pwm}).
\begin{figure}[H]
\centering
\begin{tikzpicture}[x=0.55cm,y=0.7cm,font=\scriptsize]
\draw[-{Stealth}] (0,0) -- (23,0) node[right]{$t$ (ms)};
\draw[-{Stealth}] (0,0) -- (0,1.8) node[left]{PWM};
\draw[thick,accent] (0,0) -- (0,1.2) -- (1.45,1.2) -- (1.45,0) -- (20,0) -- (20,1.2) -- (21.45,1.2) -- (21.45,0) -- (22.5,0);
\draw[{Stealth}-{Stealth}] (0,-0.5) -- node[below]{chu kỳ $T = 20$\,ms (50\,Hz)} (20,-0.5);
\draw[{Stealth}-{Stealth}] (0,1.5) -- node[above]{$p$} (1.45,1.5);
\node[align=left,anchor=west] at (3,1.0) {$p = 500$\,µs $\to 0°$\quad $p = 1450$\,µs $\to 90°$\quad $p = 2400$\,µs $\to 180°$};
\end{tikzpicture}
\caption{Dạng xung điều khiển servo SG90}\label{fig:pwm}
\end{figure}
\begin{equation}
  p(\theta) = 500 + \frac{1900\,\theta}{180}\ \text{µs}, \qquad
  D = \left\lfloor \frac{p \cdot (2^{14}-1)}{20\,000} \right\rfloor
\end{equation}
Với $\theta=90°$: $p=1450$\,µs, $D=1187$; với $\theta=0°$: $p=500$\,µs, $D=409$.

\begin{lydo}[Tại sao 50\,Hz, 14 bit, LEDC]
\textbf{50\,Hz} là tần số chuẩn của servo RC analog. \textbf{14 bit} là độ phân giải
lớn nhất của LEDC trên ESP32-C3: mỗi bước $20\,000/16\,384 \approx 1{,}22$\,µs $\approx 0{,}12°$,
mịn hơn nhiều vùng chết $\approx$1° của SG90 nên servo không rung. \textbf{LEDC} là PWM
phần cứng: xung phát liên tục mà CPU không can thiệp, servo giữ góc ổn định dù
Wi-Fi hay màn hình đang bận. Dải 500–2400\,µs là giá trị đo thực tế của SG90 (danh
nghĩa 1000–2000\,µs không đủ 180°).
\end{lydo}

\begin{lydo}[Tại sao dùng hẹn giờ một lần (\code{esp\_timer}) thay vì tác vụ hỏi vòng]
Phiên bản đầu dùng một tác vụ thức dậy 20 lần/giây để kiểm tra hạn đóng. Phiên bản
hiện tại tạo một \code{esp\_timer} một lần cho mỗi cổng: hàm đóng được gọi \emph{đúng}
lúc hết 3\,s (sai số micro giây thay vì 50\,ms), CPU không bị đánh thức vô ích và tiết
kiệm 3\,KB stack. Vì một lệnh khác có thể ``đua'' với hàm đóng, hàm đóng kiểm tra lại
hạn đóng dưới mutex; lệnh tay từ web đặt hạn về 0 và dừng hẹn giờ — \textbf{lệnh của
người vận hành luôn thắng}.
\end{lydo}

Hai cổng được mô tả bằng cùng một cấu trúc (kênh LEDC, chân, hướng lắp, trạng thái,
hạn đóng, hẹn giờ) nên một hàm xử lý cho cả hai. Thanh chắn ra lắp đối xứng nên góc
mở/đóng đảo ngược (0° mở, 90° đóng) — chỉ khai báo ở một hàm. Cả hai cổng \textbf{luôn
khởi động ở vị trí đóng}, nên trạng thái không cần lưu vào flash; phiên bản trước ghi
NVS mỗi lần thanh chắn chuyển động đã được bỏ (chậm và làm mòn flash vô ích).

\subsection{Màn hình: SPI DMA, vẽ theo sự kiện, ảnh động tiền xử lý}
\begin{lydo}[Tại sao SPI 40\,MHz và DMA]
Một khung hình đầy đủ là $240\times240\times16 = 921\,600$ bit; ở 40\,MHz mất
$\approx$ 23\,ms, tối đa $\approx$ 43 khung/giây — dư cho giao diện tĩnh và màn hình cháy
8 khung/giây. 40\,MHz bằng đúng $80\,\text{MHz}/2$ (không sai số chia tần) và dưới giới
hạn ghi $\approx$ 62\,MHz của ST7789, chừa biên an toàn cho dây cắm. DMA cho phép CPU
làm việc khác trong lúc ảnh được đẩy ra.
\end{lydo}

\begin{itemize}
\item \textbf{Framebuffer RGB565} 115\,200 byte: toàn bộ giao diện được vẽ vào RAM rồi gửi
  một lần. Màu được lưu sẵn theo thứ tự byte của ST7789 (byte cao trước) vì
  ESP32-C3 là little-endian.
\item \textbf{Hàm tô hình chữ nhật} cắt biên một lần và tô từng dòng liên tục thay vì
  kiểm tra biên và đảo byte cho từng điểm ảnh.
\item \textbf{Font 5$\times$7 tự định nghĩa} cho A–Z, 0–9 (180 byte) — không cần thư viện
  đồ họa lớn như LVGL cho một giao diện cố định.
\item \textbf{Vẽ theo sự kiện}: \file{udp\_server} đánh thức tác vụ màn hình bằng thông
  báo tác vụ ngay khi có dữ liệu; nếu không có gì, tác vụ chỉ thức mỗi 1\,s (để hết hạn
  hóa đơn/thông báo) hoặc 120\,ms khi có cháy (để lửa nhấp nháy). Màn hình chỉ vẽ lại
  khi nội dung thực sự thay đổi.
\end{itemize}

\begin{table}[H]
\centering\small
\caption{Dữ liệu ảnh động sau tiền xử lý và nén LZSS}\label{tab:lzss}
\begin{tabular}{@{}lrrrr@{}}
\toprule
\textbf{Ảnh} & \textbf{Kích thước} & \textbf{Số khung} & \textbf{Dữ liệu thô} & \textbf{Sau nén} \\ \midrule
Kaoruko & 96$\times$96 & 10 & 184\,320 B & 70\,564 B (2,6$\times$) \\
Claude & 68$\times$46 & 18 & 112\,608 B & 16\,204 B (6,9$\times$) \\
\bottomrule
\end{tabular}
\end{table}
\noindent Hai bộ đệm DMA $96\times96$ của ảnh Kaoruko (36\,KB) chính là phần heap mà bắt tay HTTPS
với SePay cần. Vì ảnh động vốn dừng khi đang hiện mã QR, tác vụ \code{lcd\_animation} trả hai bộ đệm
này lại trong lúc có xe đang thanh toán và cấp phát lại sau đó; nếu lúc ấy heap chưa đủ thì chỉ bỏ
qua khung hình.

\begin{lydo}[Tại sao tiền xử lý ảnh trên máy tính]
Ảnh gốc có nền trắng cần xóa và điểm ảnh cần đảo byte. Phiên bản trước làm việc
này trên chip cho \emph{mọi khung hình}, 10 lần/giây: thuật toán loang tìm nền trắng
nối với viền (cần thêm 1,2\,KB RAM) rồi đảo byte 9\,216 điểm. Vì ảnh không bao giờ
đổi, nhóm làm một lần trên máy tính và lưu kết quả đã xử lý: firmware chỉ còn giải
nén. Kết quả được kiểm chứng giống từng điểm ảnh với cách cũ trên cả 28 khung, dữ
liệu còn nhỏ hơn (86\,768 so với 87\,977 byte). LZSS có bộ giải nén $\approx$ 40 dòng,
không cần bộ nhớ phụ và kiểm tra biên đầy đủ nên dữ liệu hỏng không gây ghi tràn.
\end{lydo}

\textbf{Hai bộ đệm DMA} (2 $\times$ 18\,432 byte) được dùng luân phiên cho ảnh động:
lệnh vẽ chỉ \emph{xếp hàng} giao dịch DMA rồi trả về, nên giải nén khung kế tiếp vào
cùng bộ đệm khi DMA chưa đọc xong sẽ làm hình bị xé. Hai tác vụ cùng vẽ lên một
panel nên dùng chung một mutex.

\subsection{Bộ kiểm tra bản tin JSON}
\begin{lydo}[Tại sao tự viết thay vì dùng thư viện cJSON]
Gói UDP là \emph{ranh giới tin cậy}: bất kỳ thiết bị nào trong mạng đều gửi được.
cJSON cấp phát heap cho từng nút (phân mảnh heap khi chạy lâu) và chấp nhận mọi JSON
hợp lệ. Bộ đọc tự viết (\file{json\_scan.c}, 81 dòng, dùng chung cho cả hai loại bản
tin) \textbf{không cấp phát động}, chỉ chấp nhận đúng lược đồ (Bảng~\ref{tab:fields})
và có \textbf{vòng lặp giới hạn} (tối đa 4 khóa) nên thời gian xử lý bị chặn trên.
Bộ đệm nhận 256 byte, đọc tối đa 255 byte và luôn kết thúc bằng \code{'\textbackslash0'}.
\end{lydo}

\subsection{Phiên gửi xe, phí, thông báo và nhật ký}
\begin{itemize}
\item \textbf{Phiên}: mảng tĩnh 5 phần tử — bằng số chỗ đỗ, vì không thể có nhiều xe hơn số chỗ.
  Thời điểm vào lấy từ \code{esp\_timer\_get\_time()} (64 bit, micro giây, không tràn trong
  hơn 290\,000 năm).
\item \textbf{Phí}: $h=\max(1,\lceil t/3600\rceil)$, phí $=5000\,h$ VNĐ, tính bằng số
  nguyên \code{(t + 3599) / 3600}, bão hòa ở \code{UINT32\_MAX} nếu tràn.
\item \textbf{Thông báo từ chối} (sai cổng, bãi đầy): giữ 5\,s; lần quẹt hợp lệ tiếp
  theo xóa thông báo ngay.
\item \textbf{Nhật ký}: bộ đệm vòng 10 sự kiện ($\approx$ 520 byte), đọc ra theo thứ tự mới
  nhất trước; không dùng heap.
\end{itemize}

\begin{lydo}[Tại sao chọn các thông số này]
\textbf{5.000\,VNĐ/giờ, làm tròn lên, tối thiểu 1 giờ} là cách tính phổ biến của bãi
xe, dễ giải thích. \textbf{Hóa đơn 60\,s} đủ để tài xế đọc và trả tiền; \textbf{thông
báo 5\,s} đủ đọc mà không che hóa đơn quá lâu. \textbf{Thanh chắn mở 3\,s} đủ cho xe mô
hình đi qua nhưng không đủ để xe sau ``đi ké''. \textbf{Nhật ký 10 sự kiện} vừa một màn
hình điện thoại và tốn dưới 1\,KB RAM.
\end{lydo}

\subsection{Thanh toán VietQR qua SePay}\label{sec:vietqr}
Khi thanh toán QR được bật, quẹt thẻ ở cổng ra \emph{không} mở cổng ngay: máy chủ chốt
phí tại thời điểm quẹt, tạo một \textbf{mã thanh toán} và vẽ mã VietQR lên TFT. Mọi ứng
dụng ngân hàng Việt Nam quét được mã này, số tiền và nội dung chuyển khoản đã điền sẵn.
Thanh chắn ra chỉ mở khi máy chủ thấy giao dịch khớp (qua SePay) hoặc người vận hành bấm
\code{Confirm paid} trên trang web (Hình~\ref{fig:seq-pay}).

\paragraph{Chuỗi VietQR.} Hàm \code{vietqr\_payload()} (\file{parking/payment.c}) dựng
chuỗi theo chuẩn VietQR của NAPAS, vốn là định dạng EMVCo: mỗi trường gồm mã 2 chữ số,
độ dài 2 chữ số và giá trị; trường 38 và 62 chứa các trường con theo cùng cách
(Bảng~\ref{tab:vietqr}).
\begin{table}[H]
\centering\small
\caption{Cấu trúc chuỗi VietQR do máy chủ sinh}\label{tab:vietqr}
\begin{tabularx}{\linewidth}{@{}llX@{}}
\toprule
\textbf{Mã} & \textbf{Trường} & \textbf{Giá trị} \\ \midrule
00 & Phiên bản định dạng & \code{01} \\
01 & Kiểu mã & \code{12}: mã động, dùng một lần cho một số tiền \\
38 & Thông tin người nhận & 00: \code{A000000727} (định danh NAPAS); 01: \{00: mã BIN ngân hàng \code{970418} (BIDV), 01: số tài khoản\}; 02: \code{QRIBFTTA} (chuyển nhanh tới tài khoản) \\
53 & Tiền tệ & \code{704} (VNĐ) \\
54 & Số tiền & Phí đã chốt, ví dụ \code{10000} \\
58 & Quốc gia & \code{VN} \\
62 & Dữ liệu bổ sung & 08: nội dung chuyển khoản = mã thanh toán, ví dụ \code{PK7XK2QM} \\
63 & Mã kiểm tra & CRC16-CCITT (đa thức \code{0x1021}, giá trị đầu \code{0xFFFF}) tính trên toàn bộ chuỗi kể cả \code{6304}, ghi 4 chữ số hex in hoa \\
\bottomrule
\end{tabularx}
\end{table}
Hàm được kiểm chứng từng byte với một mã VietQR mẫu đã công bố (kết thúc bằng
\code{6304AB76}) và với giá trị kiểm tra chuẩn của CRC16-CCITT (\code{0x29B1} cho chuỗi
\code{"123456789"}). BIN phải đủ 6 chữ số và số tài khoản chỉ gồm chữ/số; nếu sai hoặc để
trống, máy chủ không vẽ mã QR (tránh một mã không quét được) mà hiện \code{PAY AT BOOTH}; cổng ra
vẫn chỉ mở khi người vận hành bấm \emph{Confirm paid}.

\paragraph{Mã thanh toán và lượt thanh toán.} Mã gồm \code{PK} và 6 ký tự lấy từ 31 chữ in hoa
và chữ số, đã bỏ \code{I}, \code{L}, \code{O}, \code{0}, \code{1} vì dễ nhầm khi người lái
phải gõ tay. Mỗi lúc chỉ có \textbf{một lượt thanh
toán} ở cổng ra, vì chỉ có một làn ra và một màn hình:
\begin{itemize}[nosep]
  \item Quẹt lại cùng thẻ: hiện lại đúng mã QR cũ (phí và mã không đổi), đếm lại 180\,s.
  \item Thẻ khác quẹt ra khi lượt trước chưa hết hạn: \code{RFID\_SCAN\_PAY\_BUSY},
        màn hình QR vẫn giữ cho xe thứ nhất, dòng cuối báo \code{NEXT CAR: WAIT, ONE IS PAYING}; nhật ký \code{WAIT · PAYING}.
  \item Quá 180\,s chưa có tiền: TFT báo \code{PAY AT BOOTH / NO TRANSFER YET - SEE
        OPERATOR} trong 30\,s; xe khác được bắt đầu thanh toán, nhưng chuyển khoản đến
        muộn vẫn hoàn tất lượt cũ nếu chưa bị thay thế.
  \item Hoàn tất (giao dịch khớp hoặc người vận hành xác nhận) đi qua đúng đường ra bình
        thường: hóa đơn 60\,s, nhật ký \code{OUT}, doanh thu được cộng tại thời điểm thanh
        toán, thanh chắn ra mở 3\,s. Một xe không thể ra hai lần.
\end{itemize}
Trang web nhận lượt đang chờ qua trường \code{payment} của \code{/api/rfid} và hiện băng
rôn vàng \code{WAITING FOR PAYMENT} (thẻ, phí, mã) cùng nút \code{Confirm paid} trong mục 04.

\paragraph{Màn hình thanh toán.} Mã QR được sinh bằng thành phần \code{espressif/qrcode}
(thư viện qrcodegen của Nayuki), mức sửa lỗi M, phiên bản tối đa 10, mỗi mô-đun
$\approx$ 3 điểm ảnh, viền trắng 2 mô-đun. Bên dưới là \code{SCAN TO PAY}, phí (VNĐ),
\code{NOTE <mã>  ID <thẻ>} và \code{BIDV <số tài khoản> <tên>}. Thứ tự ưu tiên hiển thị:
màn hình cháy > màn hình thanh toán > màn hình thường. Việc sinh QR không cần mạng.

\begin{figure}[H]
\centering
\mmd[0.8][0.62]{seq_pay}
\caption{Biểu đồ tuần tự thanh toán VietQR ở cổng ra, xác nhận tự động qua SePay}\label{fig:seq-pay}
\end{figure}

\paragraph{Kiểm tra chuyển khoản qua SePay.} SePay đọc biến động số dư của tài khoản
nhận và cho truy vấn qua API. Tác vụ \code{sepay} (\file{network/sepay\_client.c}, stack
8\,KB, ưu tiên 3) chỉ được tạo khi có token. Khi một lượt thanh toán bắt đầu, tác vụ gửi
\code{GET} tới \url{https://my.sepay.vn/userapi/transactions/list} với tham số
\code{limit=1} (header \code{Authorization: Bearer <token>}) để ghi lại id giao dịch mới nhất; sau đó
cứ 3\,s gọi lại với \code{since\_id=<id+1>\&limit=20}, nên chỉ giao dịch \emph{sau} lúc hiện QR
được xét. Hàm \code{sepay\_has\_payment()} tìm giao dịch có \code{transaction\_content}
chứa mã (không phân biệt hoa thường, phòng khi ngân hàng đổi kiểu chữ của nội dung) và
\code{amount\_in} $\ge$ phí. Ví MoMo và một số ví điện tử thay nội dung bằng chữ của riêng họ (đã gặp khi thử thật), nên hàm \code{sepay\_has\_transfer\_to()} còn chấp nhận giao dịch vào đúng tài khoản/VA nhận với số tiền \emph{đúng bằng} phí; vì mỗi lúc chỉ một xe thanh toán và chỉ xét giao dịch sau lúc hiện QR, cách này không nhầm xe. TLS dùng bộ chứng chỉ gốc của ESP-IDF; với
\code{CONFIG\_MBEDTLS\_DYNAMIC\_BUFFER=y}, bộ đệm TLS chỉ được cấp phát trong lúc gửi yêu
cầu — cần thiết vì RAM đã chật bên cạnh framebuffer 115\,KB. Lỗi mạng hoặc lỗi HTTP chỉ
ghi cảnh báo và thử lại ở chu kỳ sau.

\begin{table}[H]
\centering\small
\caption{Cấu hình thanh toán (\code{idf.py menuconfig → Parking QR payment (SePay)})}\label{tab:payconfig}
\begin{tabularx}{\linewidth}{@{}lX@{}}
\toprule
\textbf{Tham số} & \textbf{Ý nghĩa} \\ \midrule
\code{PARKING\_PAY\_BANK\_BIN} & Mã BIN ngân hàng nhận, mặc định \code{970418} (BIDV) \\
\code{PARKING\_PAY\_VA} & Số tài khoản nhận ghi vào mã QR. Với BIDV cá nhân phải là số \emph{tài khoản ảo (VA)} tạo trên SePay, vì SePay không nhận giao dịch chuyển thẳng vào số tài khoản chính; để trống thì không có mã QR, TFT hiện \code{PAY AT BOOTH} và cổng ra chỉ mở khi người vận hành bấm \emph{Confirm paid} \\
\code{PARKING\_PAY\_ACCOUNT\_NAME} & Tên chủ tài khoản, hiển thị trên TFT \\
\code{PARKING\_SEPAY\_API\_TOKEN} & API token SePay; để trống = người vận hành xác nhận trên trang web \\
\code{PARKING\_PAY\_TIMEOUT\_S} & Thời gian chờ chuyển khoản, mặc định 180\,s \\
\code{PARKING\_PAY\_POLL\_MS} & Chu kỳ hỏi SePay, mặc định 3000\,ms \\
\bottomrule
\end{tabularx}
\end{table}
Các giá trị này chỉ nằm trong \file{Server/sdkconfig} (bị \file{.gitignore} bỏ qua).
Token SePay có toàn quyền đọc giao dịch của tài khoản (SePay không cho giới hạn quyền
token) nên tuyệt đối không được commit.

\begin{lydo}[Tại sao hỏi vòng SePay thay vì webhook]
SePay hỗ trợ webhook (gọi tới một URL khi có tiền vào), nhưng ESP32-C3 nằm sau router
(NAT) nên không có địa chỉ mà Internet gọi tới được; dùng webhook sẽ phải thêm một máy
chủ trung gian trên Internet hoặc mở cổng trên router. Hỏi vòng chỉ cần kết nối
\emph{đi ra}, và chỉ chạy khi có xe đang thanh toán. Chu kỳ 3\,s cho độ trễ vài giây,
đủ nhanh với người đang chờ ở cổng, và thấp hơn rất nhiều so với giới hạn 2–3 yêu
cầu/giây cho mỗi token của SePay. Truy vấn \code{since\_id} giữ câu trả lời ở vài giao
dịch mới nhất nên bộ đệm nhận chỉ cần 6\,KB.
\end{lydo}

\begin{lydo}[Tại sao vẫn cho người vận hành xác nhận tay]
Thanh toán tự động phụ thuộc Internet, SePay và việc có token. Khi một trong ba thiếu
— hoặc người lái trả tiền mặt — bãi vẫn phải cho xe ra. Mã QR vẫn hiển thị vì không cần
mạng; người vận hành thấy giao dịch trong ứng dụng ngân hàng và bấm \code{Confirm paid}
(\code{POST /api/payment/confirm}, cần mã PIN quản trị). Lệnh này cũng dùng được cho xe
chưa quẹt ra (tính phí đến thời điểm xác nhận), nên một đường xử lý chung phục vụ cả
chuyển khoản lẫn tiền mặt.
\end{lydo}

\subsection{Sensor Client: lọc nhiễu và gửi theo thay đổi}
\begin{figure}[H]
\centering
\begin{tikzpicture}[x=0.9cm,y=0.6cm,font=\scriptsize]
\node[anchor=east] at (-0.2,2.5) {Tín hiệu thô};
\draw[thick] (0,2) -- (2,2) -- (2,3) -- (2.3,3) -- (2.3,2) -- (2.5,2) -- (2.5,3) -- (2.7,3) -- (2.7,2) -- (2.9,2) -- (2.9,3) -- (11,3);
\node[anchor=east] at (-0.2,0.5) {Trạng thái};
\draw[thick,accent] (0,0) -- (3.9,0) -- (3.9,1) -- (11,1);
\draw[dashed,gray] (2.9,3.3) -- (2.9,-0.4);
\draw[dashed,gray] (3.9,3.3) -- (3.9,-0.4);
\draw[{Stealth}-{Stealth}] (2.9,-0.6) -- node[below]{ổn định 50\,ms} (3.9,-0.6);
\node[align=center] at (2.45,3.8) {nhiễu khi xe\\đang lăn vào};
\node[draw,fill=boxB,align=center] at (7.5,-0.9) {gửi \code{parking\_update} ngay};
\draw[-{Stealth},thick] (3.9,1) to[bend right=10] (6.2,-0.6);
\end{tikzpicture}
\caption{Chống rung: chỉ chấp nhận mức mới khi giữ ổn định 50\,ms, rồi gửi ngay}\label{fig:debounce}
\end{figure}

\begin{lydo}[Tại sao lọc 50\,ms và gửi theo thay đổi kèm nhắc lại 1\,s]
Mép xe và bánh xe làm tín hiệu IR bật/tắt liên tục trong vài chục mili giây khi xe
đang vào chỗ. Nhóm tinh chỉnh thực nghiệm từ 250\,ms xuống \textbf{50\,ms}: đủ lọc nhiễu
mà vẫn cho cảm giác tức thời. Trước đây trạng thái chỉ được gửi theo chu kỳ 1\,s nên thay đổi
có thể trễ đến 1\,s; nay \textbf{gửi ngay khi có thay đổi}, đồng thời vẫn \textbf{nhắc
lại mỗi giây} để một gói bị mất được sửa ở lần sau. Mỗi cảm biến có bộ đếm riêng.
\end{lydo}

Mực nước được quy đổi tuyến tính từ ADC 10 bit:
\begin{equation}
  h_{\text{mm}} =
  \begin{cases}
  0, & raw \le 20 \\[2pt]
  \dfrac{\min(raw, 700) - 20}{700 - 20} \cdot 40, & raw > 20
  \end{cases}
\end{equation}
trong đó 20 là mức nhiễu khi khô, 700 là giá trị khi ngập hoàn toàn và 40\,mm là
chiều dài vùng đo; ADC được đọc mỗi 250\,ms vì đọc quá dày ảnh hưởng Wi-Fi trên
ESP8266 trong khi mực nước thay đổi chậm.

\subsection{RFID Client: hai đầu đọc trên một bus SPI}
\begin{lydo}[Tại sao bật ăng-ten luân phiên, chờ 4\,ms, gain 48\,dB, chống lặp 2\,s]
Hai đầu đọc 13,56\,MHz đặt gần nhau sẽ \textbf{nhiễu chéo}; chỉ bật một ăng-ten tại
một thời điểm loại bỏ hoàn toàn vấn đề. Thẻ thụ động cần vài mili giây để nhận đủ năng
lượng: \textbf{4\,ms} đủ cho thẻ MIFARE mà chu kỳ quét vẫn ngắn. \textbf{Gain 48\,dB} bù
tầm đọc khi đầu đọc nằm sau vỏ mô hình. Thẻ để yên sẽ bị đọc liên tục; \textbf{chống
lặp 2\,s} áp dụng cho cả hai đầu đọc, tránh một lần quẹt bị gửi nhiều lần làm ``vào''
biến thành ``ra''. UID được đổi sang chuỗi hex bằng \textbf{bộ đệm cố định} thay vì
\code{String} của Arduino: không cấp phát heap mỗi lần quẹt, nên heap 80\,KB của
ESP8266 không bị phân mảnh khi chạy lâu. DHT11 đọc mỗi \textbf{5\,s} vì cảm biến cần ít
nhất 1\,s giữa hai lần đọc và nhiệt độ thay đổi chậm.
\end{lydo}

\subsection{Trang web}
\begin{lydo}[Tại sao hỏi vòng 1–2\,s thay vì WebSocket]
Hỏi vòng không giữ trạng thái kết nối trên máy chủ: mỗi yêu cầu độc lập, xử lý xong
giải phóng ngay — phù hợp với giới hạn socket và RAM, và đơn giản khi nhiều trình duyệt
cùng xem. Với dữ liệu thay đổi theo giây, độ trễ thêm tối đa 1\,s là chấp nhận được.
Trang web là HTML/CSS/JavaScript thuần, nhúng trong firmware bằng
\code{EMBED\_TXTFILES}: không cần phân vùng tệp, không phụ thuộc CDN (trang vẫn mở được
khi router mất Internet).
\end{lydo}


\section{Triển khai trên phần cứng thực tế}\label{chap:huongdanchay}

Mục này hướng dẫn từ lúc lấy mã nguồn đến lúc hệ thống chạy hoàn chỉnh, kèm kết
quả mong đợi ở mỗi bước để kiểm tra. Môi trường cần được cài đặt trước theo
Mục~\ref{chap:caidat}. Có hai cách chạy:
\begin{itemize}
  \item \textbf{Đầy đủ phần cứng}: Bước~1 đến Bước~9.
  \item \textbf{Không cần client} (chỉ có bo máy chủ, dùng máy tính giả lập hai client):
        Bước~1, 2, 4 rồi Bước~10.
\end{itemize}
Bộ kiểm thử tự động (không cần phần cứng nào) nằm ở Bước~11.

\begin{table}[H]
\centering\small
\caption{Tên cổng nối tiếp theo hệ điều hành}\label{tab:ports}
\begin{tabularx}{\linewidth}{@{}lXX@{}}
\toprule
\textbf{Hệ điều hành} & \textbf{ESP32-C3 SuperMini (USB tích hợp)} & \textbf{ESP8266 NodeMCU (CH340/CP2102)} \\ \midrule
Windows & \code{COM3}, \code{COM4}, \ldots{} (xem Device Manager) & \code{COM5}, \ldots \\
macOS & \code{/dev/cu.usbmodem*} & \code{/dev/cu.usbserial-*} hoặc \code{/dev/cu.wchusbserial*} \\
Linux & \code{/dev/ttyACM0} & \code{/dev/ttyUSB0} \\
\bottomrule
\end{tabularx}
\end{table}

\buoc{Lấy mã nguồn}
\begin{lstlisting}[style=shell]
git clone https://github.com/thinhdinh25/INT3108_embedded-Smart-Parking-Project.git
cd INT3108_embedded-Smart-Parking-Project
\end{lstlisting}
Hoặc tải tệp ZIP từ trang GitHub của dự án và giải nén.

\buoc{Làm quen cấu trúc thư mục}
\begin{lstlisting}[language={},numbers=none,caption={Cấu trúc kho mã nguồn}]
Server/              Firmware may chu ESP32-C3 (ESP-IDF)
  main/              Ma nguon: barrier, buzzer, display, network, parking
  sdkconfig.defaults Cau hinh mac dinh: -O2, phan vung ung dung 1,5 MB
SensorClient/        Sketch Arduino cho Sensor Client (ESP8266)
RFIDClient/          Sketch Arduino cho RFID Client (ESP8266)
  secrets.example.h  Ten/mat khau Wi-Fi chung mau -> chep thanh secrets.h
examples/            Cong cu gia lap client bang Python
tests/               Bo kiem thu tu dong tren may tinh
docs/                Ghi chu dau day
report/              Bao cao nay (LaTeX + PDF)
\end{lstlisting}

\buoc{Lắp mạch}
Đấu dây theo Bảng~\ref{tab:pin-srv} (máy chủ) và Bảng~\ref{tab:pin-clients} (hai client);
bố trí thiết bị trên mô hình như Hình~\ref{fig:model} và Bảng~\ref{tab:placement}.
Kiểm tra trước khi cấp nguồn:
\begin{itemize}
  \item Servo dùng nguồn 5\,V riêng, \textbf{GND nối chung} với ESP32-C3; không cấp
        nguồn cho servo từ chân 3V3.
  \item Màn hình ST7789 dùng 3V3; RC522 dùng 3V3 (không dùng 5\,V).
  \item Cắm đúng cực jack DC (nếu cắm ngược, diode sẽ chặn nhưng khi đó hệ thống không lên nguồn);
        tụ lọc đặt sát chỗ servo lấy nguồn.
  \item Chỉnh biến trở từng module IR khi đã đặt xe mẫu vào ô: đèn báo sáng khi có xe,
        tắt khi ô trống.
  \item Module còi không được kéo GPIO9 xuống mức thấp khi cấp nguồn (nếu không, bo
        sẽ vào chế độ nạp firmware).
\end{itemize}

\buoc{Biên dịch và nạp máy chủ ESP32-C3}
Mở terminal ESP-IDF (Mục~\ref{chap:caidat}), cắm ESP32-C3 qua USB:
\begin{lstlisting}[style=shell]
cd Server
idf.py set-target esp32c3        # chi lan dau; tao sdkconfig tu sdkconfig.defaults
idf.py build                     # lan dau tai them espressif/mdns, espressif/qrcode
idf.py -p <CONG> flash monitor   # vi du: -p COM4 hoac -p /dev/cu.usbmodem101
\end{lstlisting}
\textbf{Kết quả mong đợi} sau \code{idf.py build}:
\begin{lstlisting}[language={},numbers=none]
INT3108_Smart_Parking_Project.bin binary size 0x12a9c0 bytes. Smallest app
partition is 0x177000 bytes. 0x4c640 bytes (20%) free.
Project build complete.
\end{lstlisting}
\textbf{Kết quả mong đợi} trên cửa sổ monitor (thoát bằng \code{Ctrl+]}):
\begin{lstlisting}[language={},numbers=none]
I (...) LCD_DISPLAY: ST7789 display started at 240x240 with 10-frame Kaoruko animation
I (...) WIFI_STA: Wi-Fi station started, connecting to SSID=Parking...
I (...) HTTP_SERVER: HTTP server listening on port 80
I (...) UDP_SERVER: UDP server listening on port 3333
I (...) SEPAY: QR payment on, auto-checked with SePay (timeout 180 s)
I (...) SERVER: Joining shared Wi-Fi: ESP32-Parking
I (...) SERVER: Web interface: http://parking.local
I (...) WIFI_STA: Got IP: 192.168.1.23
I (...) SERVER: mDNS: parking.local -> 192.168.1.23
\end{lstlisting}
\code{Joining shared Wi-Fi} in tên mạng Wi-Fi chung đã cấu hình; \code{Got IP} là địa chỉ
do router cấp (ở đây \code{192.168.1.23}, tùy router). Dòng \code{SEPAY} cho biết chế độ
thanh toán: \code{No QR} nếu chưa khai báo tài khoản nhận (cổng ra chờ người vận hành xác nhận),
\code{confirmed by the operator} nếu có tài khoản nhưng chưa có token SePay. Màn hình TFT
hiện giao diện \code{PARKSENSE} với các giá trị \code{--.-}; hai thanh chắn về vị trí
đóng. Nếu chưa thấy \code{Got IP}, kiểm tra router đã bật và SSID/mật khẩu đúng.

\begin{lydo}[Đổi mật khẩu Wi-Fi và mã PIN trước khi triển khai]
Mã nguồn chỉ chứa giá trị \emph{demo}. Giá trị thật không nằm trong Git:
\begin{itemize}[nosep]
  \item Máy chủ: \code{idf.py menuconfig → Parking network and security}, đặt SSID, mật
        khẩu của mạng Wi-Fi chung và mã PIN quản trị. Giá trị
        được lưu trong \file{Server/sdkconfig}, tệp này bị \file{.gitignore} bỏ qua.
  \item Thanh toán QR: \code{idf.py menuconfig → Parking QR payment (SePay)}, đặt số tài
        khoản nhận, tên chủ tài khoản và token SePay (Bảng~\ref{tab:payconfig}). Để trống
        số tài khoản thì không có mã QR, người vận hành thu tiền và bấm \emph{Confirm paid}. Token SePay có toàn quyền đọc giao dịch của tài
        khoản, nên \textbf{chỉ} được nhập ở đây (\file{sdkconfig}), không bao giờ ghi vào
        \file{sdkconfig.defaults} hay mã nguồn.
  \item Hai client: sao chép \file{secrets.example.h} thành \file{secrets.h} trong thư mục
        của sketch rồi sửa; \file{secrets.h} cũng bị bỏ qua khi commit. Không có
        \file{secrets.h} thì sketch dùng giá trị demo.
  \item SSID và mật khẩu Wi-Fi phải \textbf{giống nhau} trên cả ba bo (cùng một mạng chung).
\end{itemize}
\end{lydo}

\begin{lydo}[Nâng cấp từ bản firmware cũ]
Bản này thêm thống kê doanh thu vào dữ liệu lưu trong NVS nên định dạng lưu đổi
(\code{RFID\_STORE\_VERSION} 3). Lần khởi động đầu sau khi nạp, máy chủ bỏ qua bản lưu
cũ: các xe đang gửi phải quẹt vào lại. Các lần sau dữ liệu được giữ bình thường.
\end{lydo}

\begin{lydo}[Lưu ý khi đã có cấu hình cũ]
\file{sdkconfig.defaults} chỉ được áp dụng khi chưa có tệp \file{Server/sdkconfig}.
Nếu máy đã từng biên dịch bản cũ, hãy chạy lại \code{idf.py set-target esp32c3} (tạo lại
cấu hình) hoặc chọn trong \code{idf.py menuconfig}: \emph{Compiler options → Optimize for
performance} và \emph{Partition Table → Single factory app (large), no OTA}. Lần nạp
đầu sau khi đổi bảng phân vùng phải dùng \code{idf.py flash} (nạp cả bảng phân vùng).
\end{lydo}

\buoc{Nạp RFID Client}
\begin{enumerate}
  \item Mở Arduino IDE, \emph{File → Open} tệp \file{RFIDClient/RFIDClient.ino}.
  \item \emph{Tools → Board}: \textbf{NodeMCU 1.0 (ESP-12E Module)}; \emph{Port}: cổng
        của bo (Bảng~\ref{tab:ports}); \emph{Upload Speed}: 115200.
  \item Bấm \textbf{Upload} (mũi tên). Chờ \code{Done uploading}.
  \item Mở \emph{Serial Monitor}, tốc độ \textbf{115200}.
\end{enumerate}
Trước khi nạp, sao chép \file{RFIDClient/secrets.example.h} thành \file{secrets.h} và
điền SSID/mật khẩu mạng Wi-Fi chung (giống máy chủ). \textbf{Kết quả mong đợi}: dòng
\code{Đang kết nối Wi-Fi <SSID>}, sau đó \code{Đã kết nối, IP: 192.168.1.x}, phiên bản
firmware của hai đầu đọc và \code{Hệ thống đã sẵn sàng!}. RFID Client chỉ là một trạm
trong mạng chung; nó gửi bản tin tới \code{255.255.255.255:3333} nên không cần biết IP
máy chủ.

\buoc{Nạp Sensor Client}
Làm tương tự Bước~5 với tệp \file{SensorClient/SensorClient.ino} (cũng tạo
\file{secrets.h} từ \file{secrets.example.h}). \textbf{Kết quả mong đợi}
trên Serial Monitor: \code{Connected, client IP: 192.168.1.x}, sau đó mỗi giây một dòng
\code{Sent: \{"type":"parking\_update",...\}}. Màn hình TFT hiện số chỗ trống thật.

\begin{lydo}[Nạp lại qua Wi-Fi (OTA)]
Từ lần thứ hai, khi máy tính kết nối cùng mạng Wi-Fi chung, hai client xuất hiện
trong \emph{Tools → Port} dưới dạng cổng mạng (ví dụ \code{ESP8266-RC522-DHT11}); chọn
cổng đó để nạp firmware mà không cần tháo bo khỏi mô hình.
\end{lydo}

\buoc{Bật hệ thống theo thứ tự}
Bật router/điểm phát Wi-Fi chung trước, sau đó bật ba bo theo \textbf{thứ tự bất kỳ}:
cả ba đều là trạm Wi-Fi, tự kết nối (và tự kết nối lại) vào mạng chung, và hai client
gửi broadcast nên không cần chờ máy chủ có IP.

\buoc{Mở trang web giám sát}
\begin{enumerate}
  \item Cho điện thoại/máy tính kết nối \textbf{cùng mạng Wi-Fi chung} với ba bo.
  \item Mở trình duyệt, vào \code{http://parking.local}.
  \item Lần đầu bấm một nút điều khiển thanh chắn, cần nhập mã PIN quản trị (mặc định
        \code{2468}). Trình duyệt nhớ mã PIN và từ đó tự gửi giờ thực cho máy chủ mỗi
        lần mở trang, nên giờ vào/ra trong nhật ký là giờ thật.
  \item Nếu thiết bị không mở được tên \code{.local} (một số máy Android cũ), dùng địa chỉ
        IP trong dòng \code{Web interface: http://192.168.1.x} (hoặc \code{Got IP}) trên monitor của máy chủ.
\end{enumerate}
\textbf{Kết quả mong đợi}: trang như Hình~\ref{fig:dashboard}, góc trên phải hiện
\code{ESP32 / ONLINE} và nhãn \code{LIVE / SENSOR}.

\buoc{Chạy các kịch bản demo}
Thực hiện kịch bản ở mục~\ref{sec:demo} và đối chiếu với Bảng~\ref{tab:systest}
(Chương~\ref{chap:kiemthu}).

\buoc{Chạy khi không có hai client (giả lập bằng máy tính)}
\begin{enumerate}
  \item Cho máy chủ và máy tính cùng kết nối vào mạng Wi-Fi chung (router, hoặc điểm phát
        trên điện thoại có tên/mật khẩu như đã đặt ở Bước~4); mở monitor máy chủ để
        thấy \code{Got IP: \ldots}.
  \item Giả lập Sensor Client và RFID Client. Script mặc định gửi tới
        \code{255.255.255.255} nên không cần \code{--host}; chỉ thêm \code{--host <IP>}
        (IP máy chủ) khi mạng chặn broadcast:
\end{enumerate}
\begin{lstlisting}[style=shell]
python3 examples/udp_sender.py --slots 1,0,0,1,0
python3 examples/udp_rfid_sender.py A1B2C3D4 --gate in
python3 examples/udp_rfid_sender.py A1B2C3D4 --gate in   # quet lai: WRONG GATE
python3 examples/udp_rfid_sender.py A1B2C3D4 --gate out  # ra: ma QR (hoac hoa don)
python3 examples/udp_sender.py --slots 1,0,0,1,0 --fire  # bao chay
python3 examples/udp_sender.py --slots 1,0,0,1,0         # het chay
\end{lstlisting}
\textbf{Kết quả mong đợi}: TFT và trang web thay đổi đúng như khi dùng client thật;
thanh chắn mở/đóng; nhật ký trên web có các dòng \code{IN}, \code{REJECTED · WRONG GATE},
\code{OUT}. Khi thanh toán QR bật, lệnh \code{--gate out} hiện mã QR thay vì mở cổng;
chuyển khoản thật hoặc bấm \code{Confirm paid} trên web để xe ra. Trên Windows dùng
\code{python} thay cho \code{python3}.

Script RFID đặt \code{seq} ngẫu nhiên cho mỗi lần chạy, nên có thể chạy song song với
RFID Client thật.

\buoc{Chạy bộ kiểm thử tự động}
Không cần ESP-IDF hay Arduino, chỉ cần trình biên dịch C/C++:
\begin{lstlisting}[style=shell]
make -C tests
\end{lstlisting}
\textbf{Kết quả mong đợi}:
\begin{lstlisting}[language={},numbers=none]
payment                  30/30 checks passed
json_parsers             78/78 checks passed
rfid_session             149/149 checks passed
gif_frames               89/89 checks passed
sensor_client            15/15 checks passed
rfid_client              23/23 checks passed
ALL TESTS PASSED
\end{lstlisting}

\subsection{Xử lý sự cố thường gặp}
\begin{longtable}{@{}p{4.6cm}p{9.8cm}@{}}
\caption{Lỗi thường gặp và cách khắc phục}\label{tab:troubleshoot}\\
\toprule \textbf{Hiện tượng} & \textbf{Nguyên nhân và cách khắc phục} \\ \midrule \endfirsthead
\toprule \textbf{Hiện tượng} & \textbf{Nguyên nhân và cách khắc phục} \\ \midrule \endhead
\code{idf.py: command not found} & Chưa mở terminal ESP-IDF; mở shortcut/script kích hoạt do EIM tạo. \\
\code{Failed to connect to ESP32-C3} & Sai cổng; hoặc giữ nút BOOT, bấm RESET, thả BOOT rồi nạp lại. \\
Biên dịch báo phân vùng đầy & Đang dùng \file{sdkconfig} cũ (1\,MB): làm theo hộp ``Lưu ý khi đã có cấu hình cũ'' ở Bước~4. \\
Monitor lặp \code{Disconnected ... retrying} & Router/điểm phát chưa bật, ngoài tầm phủ sóng, hoặc sai SSID/mật khẩu trong menuconfig; kiểm tra lại Bước~4. \\
Máy chủ nhận được bản tin nhưng client không thấy ACK, hoặc máy chủ không có dữ liệu dù client báo đã gửi & Router/điểm phát đang bật chế độ cô lập thiết bị (\emph{AP/client isolation}), chặn gói giữa các thiết bị trong mạng: tắt chế độ này trong trang quản trị router. Mạng khách (guest) thường bật sẵn chế độ này. \\
Mã QR hiện nhưng chuyển khoản xong cổng không mở & Chưa có token SePay hoặc router mất Internet (xem dòng \code{SEPAY} trên monitor); kiểm tra nội dung chuyển khoản còn nguyên mã \code{PK\ldots}, số tiền đủ; người vận hành có thể bấm \code{Confirm paid}. \\
Bấm nút trên web không có tác dụng & Sai mã PIN: xóa dữ liệu trang (hoặc mở tab ẩn danh) rồi nhập lại. \\
Arduino không thấy cổng ESP8266 & Cài driver CH340/CP2102; thử cáp USB khác (cáp chỉ sạc không truyền được dữ liệu). \\
Thiếu \code{MFRC522.h} hoặc \code{DHT.h} & Cài thư viện như Mục~\ref{chap:caidat}. \\
\code{parking.local} không mở & Thiết bị không hỗ trợ mDNS: dùng IP trong log máy chủ; kiểm tra điện thoại đã kết nối đúng mạng Wi-Fi chung với máy chủ. \\
Trang web báo \code{OFFLINE} & Máy chủ mất điện hoặc mất Wi-Fi; xem monitor máy chủ. \\
Servo rung, bo tự khởi động lại & Cấp nguồn servo từ 3V3 hoặc thiếu GND chung; dùng nguồn 5\,V riêng. \\
Màn hình trắng hoặc lệch hình & Kiểm tra dây SPI; chỉnh chân/độ lệch trong \code{idf.py menuconfig → Parking TFT display}. \\
Ô đỗ báo ngược (trống thành có xe) & Chỉnh biến trở trên module IR; kiểm tra thứ tự chân S1–S5. \\
\bottomrule
\end{longtable}
